\documentclass[11pt]{article}
\usepackage[a4paper,margin=27mm]{geometry}
\usepackage{amsmath,amssymb,amsthm,mathtools}
\usepackage{booktabs}
\usepackage{enumitem}
\usepackage{graphicx}
\usepackage[final]{microtype}
\usepackage{xcolor}
\usepackage[colorlinks=true,linkcolor=blue!55!black,citecolor=blue!55!black,urlcolor=blue!55!black]{hyperref}

\newtheorem{theorem}{Theorem}
\newtheorem{lemma}[theorem]{Lemma}
\newtheorem{corollary}[theorem]{Corollary}
\newcommand{\N}{\mathbb N}
\newcommand{\Z}{\mathbb Z}
\newcommand{\T}{\mathbb T}
\newcommand{\CalC}{\mathcal C}
\newcommand{\CalS}{\mathcal S}
\newcommand{\vtwo}{\nu_2}
\newcommand{\op}{\operatorname{op}}

\title{The exact Riesz bounds of the ReLU tent-function basis\\
\large A full solution of the optimal-bound problem in arXiv:2303.00076}
\author{Automated mathematics run \texttt{fa\_banach\_001}\\Lane 14, agent \texttt{agent\_lane\_14}}
\date{17 August 2026}

\begin{document}
\maketitle

\begin{center}
\fbox{\parbox{0.91\textwidth}{\textbf{Status: claimed full solution; likely valid, pending human review.}
The optimal uniform Gram spectral interval is
\(\boxed{[2/3,\,3/2]}\).  Thus the published lower estimate \(1/2\)
improves to the sharp value \(2/3\), while the upper estimate \(3/2\) is
already sharp.  The same exact constants hold for the normalized
multivariate system in every dimension.}}
\end{center}

\section{Source and open problem}

The source is Cornelia Schneider and Jan Vyb\'iral, \emph{A multivariate
Riesz basis of ReLU neural networks}, arXiv:2303.00076, published in
\emph{Applied and Computational Harmonic Analysis} \textbf{68} (2024),
article 101605, DOI \href{https://doi.org/10.1016/j.acha.2023.101605}{10.1016/j.acha.2023.101605}
\cite{SV}.  The question occurs in Remark 2.3(3) on PDF page 5, immediately
after Theorem 2.2.

\begin{figure}[ht]
  \centering
  \includegraphics[width=\textwidth]{figures/open_problem_crop.png}
  \caption{Full-width crop of Theorem 2.2 and Remark 2.3 from source PDF
  page 5.  The source's plots show the \emph{unnormalized} cosine Gram
  eigenvalues; see Corollary~\ref{cor:unnormalized}.}
  \label{fig:source-question}
\end{figure}

Writing
\[
 \overline{\mathcal R}_1^N
 =\{1\}\cup\{\sqrt3\,\CalC_k,\sqrt3\,\CalS_k:1\leq k\leq N\},
\]
the source proves that its Gram spectrum lies in \([1/2,3/2]\), uniformly in
\(N\), and asks whether those bounds are optimal.  It conjectures that there
is room for improvement \cite[Remark 2.3(3)]{SV}.

\section{Proof intuition}

The Gram matrix is not merely sparse; it is multiplicative.  Indices having
different powers of two are orthogonal.  Within one fixed \(2\)-adic block,
remove the common power of two and write the remaining odd integers by their
prime-exponent vectors.  The normalized cosine kernel then factors prime by
prime:
\[
 \frac{\gcd(m,n)^4}{m^2n^2}
 =\prod_{p\ \mathrm{odd}}p^{-2|v_p(m)-v_p(n)|}.
\]
For one prime this is the Toeplitz kernel \(r^{|a-b|}\), whose Fourier symbol
has exact range
\([\,(1-r)/(1+r),(1+r)/(1-r)\,]\).  Tensoring over the odd primes and
evaluating the Euler products gives \([2/3,3/2]\).

The inequalities alone would not prove optimality.  To reach both endpoints,
take finite rectangular grids of prime exponents.  Constant coefficients
concentrate the local Fourier mass at angle \(0\); alternating coefficients
concentrate it at angle \(\pi\).  These are genuine finite coefficient
vectors, hence occur inside sufficiently large initial Gram matrices.  The
sine block is a diagonal sign conjugate of the cosine block.  Finally, in
several variables, distinct primitive lattice rays are orthogonal and every
ray carries exactly the same univariate block.

\section{Statement of the result}

Let \(G_N\) be the Gram matrix of the ordered normalized system
\(\overline{\mathcal R}_1^N\).  Thus its diagonal entries are one.  The source
uses Riesz constants in the squared-norm convention
\[
 A\sum_j|a_j|^2\leq \Bigl\|\sum_j a_j u_j\Bigr\|_2^2
 \leq B\sum_j|a_j|^2.
\]

\begin{theorem}[Exact uniform constants]\label{thm:main}
For every \(N\geq1\),
\[
 \sigma(G_N)\subseteq[2/3,3/2].
\]
Moreover,
\[
 \inf_{N\geq1}\lambda_{\min}(G_N)=\frac23,
 \qquad
 \sup_{N\geq1}\lambda_{\max}(G_N)=\frac32.
\]
Consequently the optimal Riesz bounds of the infinite normalized univariate
system are exactly \(A=2/3\) and \(B=3/2\).

For every \(d\geq1\), the normalized multivariate system of \cite{SV},
\[
 \{1\}\cup
 \{\sqrt3\,\CalC(\alpha\cdot x):\alpha\in\Z^d,\ \alpha\mathbin{\triangleright}0\}
 \cup
 \{\sqrt3\,\CalS(\alpha\cdot x):\alpha\in\Z^d,\ \alpha\mathbin{\triangleright}0\},
\]
has the same exact optimal Riesz bounds \(2/3\) and \(3/2\), independently of
\(d\).  Here \(\alpha\mathbin{\triangleright}0\) means that the first nonzero
coordinate of \(\alpha\) is positive.
\end{theorem}

\section{Proof}

\subsection{The source Gram formulas}

Put \(C_N=(3\langle\CalC_i,\CalC_j\rangle)_{1\leq i,j\leq N}\) and define
\(S_N\) analogously.  Lemma 2.1 of \cite{SV} gives
\[
 G_N=[1]\oplus C_N\oplus S_N,
\]
and states that the cosine entry is zero when \(\vtwo(i)\neq\vtwo(j)\).
When \(\vtwo(i)=\vtwo(j)=a\), write \(i=2^am\) and \(j=2^an\) with \(m,n\)
odd.  Then
\begin{equation}\label{eq:odd-kernel}
 (C_N)_{ij}
 =\frac{\gcd(m,n)^4}{m^2n^2}
 =\prod_{p\ \mathrm{odd}}p^{-2|v_p(m)-v_p(n)|}.
\end{equation}
Only finitely many factors in the last product differ from one.

\subsection{The exact one-prime inequality}

\begin{lemma}\label{lem:local}
Let \(0<r<1\), and let \(A_r=(r^{|u-v|})_{u,v\geq0}\) act through its
quadratic form on finitely supported vectors.  Then
\begin{equation}\label{eq:local-bounds}
 \frac{1-r}{1+r}\,\|x\|_2^2
 \leq \sum_{u,v\geq0}r^{|u-v|}x_u\overline{x_v}
 \leq \frac{1+r}{1-r}\,\|x\|_2^2.
\end{equation}
Both constants are optimal.
\end{lemma}

\begin{proof}
The Poisson-type symbol with Fourier coefficients \(r^{|h|}\) is
\[
 f_r(\theta)=\sum_{h\in\Z}r^{|h|}e^{ih\theta}
 =\frac{1-r^2}{1-2r\cos\theta+r^2}.
\]
For finitely supported \(x\), Parseval gives
\[
 \sum_{u,v\geq0}r^{|u-v|}x_u\overline{x_v}
 =\int_0^{2\pi}f_r(\theta)
   \left|\sum_{u\geq0}x_ue^{iu\theta}\right|^2\frac{d\theta}{2\pi}.
\]
The minimum and maximum of \(f_r\) occur at \(\pi\) and \(0\), respectively,
and are precisely the two constants in \eqref{eq:local-bounds}.

For sharpness, take the normalized vectors
\(x_u=M^{-1/2}\) and \(x_u=M^{-1/2}(-1)^u\) for \(0\leq u<M\).  Their
Rayleigh quotients are
\begin{equation}\label{eq:qpm}
 q_+(r,M)=1+\frac2M\sum_{h=1}^{M-1}(M-h)r^h,
 \qquad
 q_-(r,M)=1+\frac2M\sum_{h=1}^{M-1}(M-h)(-r)^h.
\end{equation}
As \(M\to\infty\), these tend to \((1+r)/(1-r)\) and
\((1-r)/(1+r)\), respectively.
\end{proof}

\subsection{Tensoring over odd primes}

Let \(F\) be a finite set of odd primes.  On
\(\ell_2(\N_0^F)=\bigotimes_{p\in F}\ell_2(\N_0)\), define
\[
 A_F=\bigotimes_{p\in F}A_{p^{-2}}.
\]
Repeated application of Lemma~\ref{lem:local} gives
\begin{equation}\label{eq:finiteF}
 L_F I\leq A_F\leq U_F I,
 \quad
 L_F=\prod_{p\in F}\frac{1-p^{-2}}{1+p^{-2}},
 \quad U_F=L_F^{-1}.
\end{equation}

Now let \(x\) be any finitely supported coefficient vector on the odd positive
integers.  Choose \(F\) containing every prime divisor of every index in the
support.  The map \(n\mapsto(v_p(n))_{p\in F}\) identifies its quadratic form
from \eqref{eq:odd-kernel} with the restriction of the quadratic form of
\(A_F\).  Therefore \eqref{eq:finiteF} applies even though the support need
not be a rectangular exponent grid.

The Euler products are
\begin{align*}
 \prod_{p\ \mathrm{odd}}(1-p^{-2})
 &=\frac{1/\zeta(2)}{1-2^{-2}}=\frac8{\pi^2},\\
 \prod_{p\ \mathrm{odd}}(1+p^{-2})
 &=\frac{\prod_{p\ \mathrm{odd}}(1-p^{-4})}
         {\prod_{p\ \mathrm{odd}}(1-p^{-2})}
 =\frac{12}{\pi^2}.
\end{align*}
Consequently
\begin{equation}\label{eq:global-products}
 \prod_{p\ \mathrm{odd}}\frac{1-p^{-2}}{1+p^{-2}}=\frac23,
 \qquad
 \prod_{p\ \mathrm{odd}}\frac{1+p^{-2}}{1-p^{-2}}=\frac32.
\end{equation}
Since \(L_F\geq2/3\) and \(U_F\leq3/2\), every odd block is bounded between
\(2/3\) and \(3/2\).  Every fixed \(2\)-adic block is a copy of the odd block,
so
\begin{equation}\label{eq:cosine-bound}
 \frac23 I\leq C_N\leq\frac32 I
\end{equation}
for every \(N\).

\subsection{The sine block is a sign conjugate}

For \(k\geq1\), let \(k=2^{\vtwo(k)}k_{\mathrm{odd}}\) and set
\[
 \varepsilon_k=\chi_4(k_{\mathrm{odd}}),\qquad
 \chi_4(t)=\begin{cases}1,&t\equiv1\pmod4,\\-1,&t\equiv3\pmod4.
 \end{cases}
\]
Let \(D_N=\operatorname{diag}(\varepsilon_1,\ldots,\varepsilon_N)\).
We claim that
\begin{equation}\label{eq:sign-similarity}
 S_N=D_NC_ND_N.
\end{equation}
Indeed, entries from different \(2\)-adic blocks vanish on both sides.  In one
block write \(i=2^am\), \(j=2^an\), \(d=\gcd(m,n)\), and
\(m_0=m/d,n_0=n/d\).  The sign rule in Lemma 2.1 of \cite{SV} says that the
sine entry is negative exactly when \((m_0+n_0)/2\) is even.  For coprime odd
\(m_0,n_0\), this is equivalent to
\(\chi_4(m_0)\chi_4(n_0)=-1\).  Since \(\chi_4(d)^2=1\), this sign is
\(\varepsilon_i\varepsilon_j\), proving \eqref{eq:sign-similarity}.
Thus \(S_N\) and \(C_N\) have identical spectra.  This similarity was also
proved independently in the later paper \cite{KLSV}; here the displayed
argument derives it directly from the source formula.

Combining \eqref{eq:cosine-bound}, \eqref{eq:sign-similarity}, and the scalar
constant block proves \(\sigma(G_N)\subseteq[2/3,3/2]\).

\subsection{Finite witnesses prove sharpness}

Fix a finite set \(F\) of odd primes and \(M\geq1\).  Consider the finite
index set
\[
 E_{F,M}=\left\{\prod_{p\in F}p^{a_p}:0\leq a_p<M\right\}.
\]
Give these indices either constant coefficients or coefficients
\((-1)^{\sum_{p\in F}a_p}\), and normalize by \(|E_{F,M}|^{-1/2}\).
The product structure makes their cosine Rayleigh quotients exactly
\begin{equation}\label{eq:tensor-witnesses}
 Q_+(F,M)=\prod_{p\in F}q_+(p^{-2},M),
 \qquad
 Q_-(F,M)=\prod_{p\in F}q_-(p^{-2},M).
\end{equation}
For fixed \(F\), equation \eqref{eq:qpm} yields
\[
 Q_+(F,M)\longrightarrow U_F,
 \qquad Q_-(F,M)\longrightarrow L_F
 \quad(M\to\infty).
\]
Then let \(F\) increase through the odd primes and use
\eqref{eq:global-products}.  For every \(\delta>0\), this produces a finite
coefficient vector with quotient greater than \(3/2-\delta\), and another
with quotient less than \(2/3+\delta\).

Every set \(E_{F,M}\) lies inside \(\{1,\ldots,N\}\) for some finite \(N\).
Extend the witness by zeros on the other cosine indices and on the constant
and sine blocks.  The min--max principle therefore gives
\[
 \sup_N\lambda_{\max}(G_N)\geq\frac32,
 \qquad
 \inf_N\lambda_{\min}(G_N)\leq\frac23.
\]
Together with the already proved inequalities, both are equalities.  This
completes the univariate proof.

\subsection{Multivariate upgrade}

Every oriented nonzero lattice vector has a unique representation
\(\alpha=k\gamma\), where \(k\in\N\) and \(\gamma\in\Z^d\) is primitive and
oriented.  Lemma 3.1 of \cite{SV} states that functions indexed by non-collinear
vectors are orthogonal, that cosine--sine cross terms vanish, and that along
one ray \(\{k\gamma:k\geq1\}\) the normalized cosine and sine inner products
are exactly the univariate ones.  Hence, after reordering, every finite
multivariate Gram matrix is a compression of the orthogonal direct sum
\[
 [1]\ \oplus\!
 \bigoplus_{\gamma\ \mathrm{primitive,oriented}}
 \bigl(C\oplus S\bigr).
\]
The bounds \([2/3,3/2]\) therefore hold in every dimension.  The source
already proves completeness, so these are Riesz-basis bounds.  Sharpness is
realized asymptotically on the single primitive ray generated by
\((1,0,\ldots,0)\), using the finite witnesses above.  Theorem~\ref{thm:main}
is proved. \qed

\begin{corollary}[Normalization of the source plots]\label{cor:unnormalized}
For the unnormalized cosine and sine systems, whose Gram blocks are
\((\langle\CalC_i,\CalC_j\rangle)\) and
\((\langle\CalS_i,\CalS_j\rangle)\), the exact optimal uniform spectral
interval is \([2/9,1/2]\).
\end{corollary}
\begin{proof}
The normalized blocks are three times the unnormalized blocks.
\end{proof}

\section{Verification}

The proof is analytic.  The included script \texttt{code/verify\_bounds.py}
performs three independent numerical audits:
\begin{enumerate}[leftmargin=*]
\item it constructs normalized cosine Gram matrices for
\(N=16,32,64,128,256,512\) and checks all eigenvalues against
\([2/3,3/2]\);
\item it constructs \(S_N=D_NC_ND_N\) and compares the complete cosine and
sine spectra (worst mismatch: \(0.000\mathrm e{+00}\) in the run);
\item on the first eight odd primes it evaluates the length-80 finite tensor
witnesses, obtaining \(0.681311350825\) below and \(1.467744146372\) above,
moving toward \(2/3\) and \(3/2\), respectively.
\end{enumerate}
The finite spectra at \(N=512\) were
\(\lambda_{\min}=0.773698855226\) and
\(\lambda_{\max}=1.280150591375\).  The reproducible command is
\begin{verbatim}
conda run --no-capture-output -n sandbox python \
  code/verify_bounds.py
\end{verbatim}
These computations are consistency checks, not ingredients of the proof.

\section{Novelty check, limitations, and review focus}

\paragraph{Bounded novelty check.}
A bounded search through 17 August 2026 used the exact source title and
authors, citation links, arXiv:2303.00076, ``optimal Riesz bounds,'' ``Gram
spectrum,'' and the candidate constants \(2/3\) and \(3/2\).  The strongest
relevant follow-up found was Abdullo Kulbatov, Jan Lang, Cornelia Schneider,
and Jan Vyb\'iral, \emph{Bases of Lebesgue spaces formed by neural networks},
arXiv:2511.23179 \cite{KLSV}.  It proves the finite cosine/sine spectral
coincidence and improves the known lower estimate to
\[
 \left\{\frac{4\sqrt6}{\pi^2}
 \left(2-\frac{\pi^2}{8}\right)\right\}^2
 =0.5787\ldots,
\]
in both the univariate and multivariate settings, while retaining upper bound
\(3/2\).  It does not derive \(2/3\) or prove optimality.  A later paper using
the basis for Sobolev and Barron approximation likewise did not provide the
exact constants in the checked material.  No matching exact answer was found.
This is a bounded negative search, not a claim of exhaustive novelty.

\paragraph{Limitations.}
The theorem determines the best constants uniform in \(N\); it does not give a
closed formula for the extreme eigenvalues at each fixed \(N\), nor rates at
which selected initial sections approach the endpoints.  The multivariate
statement improves the constants for the precise lattice-ridge basis of
\cite{SV}; it is not a claim about arbitrary ReLU dictionaries.

\paragraph{Human-review recommendation.}
Review the identification of a finitely supported odd-index quadratic form
with a compression of the finite-prime tensor product, the order of the two
limits in the explicit sharpness witnesses, and the primitive-ray direct-sum
decomposition in the multivariate argument.  Also verify that the source's
Riesz constants use the squared-norm convention, as adopted here.

\begin{thebibliography}{9}
\bibitem{SV}
C. Schneider and J. Vyb\'iral,
\emph{A multivariate Riesz basis of ReLU neural networks},
Applied and Computational Harmonic Analysis \textbf{68} (2024), 101605;
arXiv:2303.00076.

\bibitem{KLSV}
A. Kulbatov, J. Lang, C. Schneider, and J. Vyb\'iral,
\emph{Bases of Lebesgue spaces formed by neural networks},
arXiv:2511.23179 (2025).
\end{thebibliography}

\end{document}
